% !TEX program = lualatex
\documentclass[a4paper,12pt]{article}
\usepackage{fontspec}
\usepackage{polyglossia}
\usepackage{amsmath}
\usepackage{microtype}
\usepackage{unicode-math}
\usepackage[left=15mm,right=15mm,top=15mm,bottom=20mm]{geometry}
\usepackage{graphicx}
\usepackage{caption}
% \usepackage{tabularx}
\usepackage{wrapfig}
\usepackage{mdframed}
\usepackage{setspace}

\setmainlanguage{russian}
\setotherlanguage{english}
% Переключение: \textenglish{...} или \begin{english}...\end{english}
\setmainfont{Times New Roman}
\setmonofont{Source Code Pro}
% \setmathfont{Libertinus Math}
% \setmathfont{STIX Two Math}
\setmathfont{Cambria Math}
\setlength{\parindent}{10mm}
\setlength{\parskip}{2mm}

\pretocmd{\subsection}{\clearpage}{}{}

\begin{document}

		\paragraph*{Расширение тел.} \

		\textit{Изменение размеров тела пропорционально изменению температуры.}

		Тепловое расширение материала характеризуется величиной

		\begin{equation}
			\alpha = \frac{l-l_0}{l_0(t-t_0)}
		\end{equation}

		где $t_0$ - начальная температура, $l_0$ - начальная длина, $t$ - текущая температура,
		$l$ - текущая длина.

		Эта величина называется \textit{температурным коэффициентом линейного расширения} и 
		показывает, на какую долю своей нормальной длины увеличивается длина тела при 
		нагревании на $1\text{К}$.

		\textit{Температурный коэффициент объемного расширения} характеризуется величиной

		\begin{equation}
			\beta = \frac{V-V_0}{V_0(t-t_0)}
		\end{equation}

		где $t_0$ - начальная температура, $V_0$ - начальный объем, $t$ - текущая температура,
		$V$ - текущий объем.

		При увеличении объема тел плотность их уменьшается во столько раз, во сколько
		увеличился объем.
		
		\begin{equation}
			\rho = \frac{\rho_0}{1+\beta(t-t_0)}
		\end{equation}


		\subsection*{Теплота}

		Внутренняя энергия тела зависит от его температуры, от того, является ли тело твердым, 
		жидким или газообразным, находится ли оно в мелкораздробленном состоянии или является сплошным, 
		и т. д. Если под действием внешней силы (сила, приложенная к телу извне, со стороны другого 
		тела или объекта, например рука, мотор) производится работа против сил трения, в результате чего 
		температура тела повышается или оно измельчается, расплавляется или испаряется, то внутренняя 
		энергия тела увеличивается. Если, наоборот, температура тела понижается, если оно превращается 
		из газообразного в жидкое и т. п., то внутренняя энергия тела уменьшается.

		Cумма кинетической, потенциальной и внутренней энергий всех участвующих тел не изменяется. 
		Эта сумма - полная энергия тела.

		Не всегда изменение внутренней энергии тела может происходить при совершении работы. 
		В этих случаях говорят о передаче теплоты.

		\textit{Передача теплоты - процесс, при котором внутренняя энергия одних тел уменьшается, 
		а других — соответственно увеличивается, причем механическая энергия тел не изменяется и 
		никакая работа не совершается.}

		\textit{Количество теплоты - это изменение внутренней энергии тела, которое происходит при
		теплопередачи. Единица количества теплоты - джоуль (Дж).}

		При определенном изменении температуры изменение внутренней энергии тела пропорционально
		его массе.

		Количество теплоты, которое нужно сообщить какому-либо телу, чтобы повысить его температуру
		на 1К, называется \textit{теплоемкостью} этого тела. При остывании тело отдает такое же
		количество теплоты.
		
		Теплоемкость тела пропорциональна массе тела и зависит от вещества, из которого оно
		состоит.


\end{document}